% Solutions to Chapter 4's exercises, from lectures/L04/appendix/d_solutions.md. Every exercise in
% the chapter is worked on paper, so every one is answered. The appendix's opening paragraph, which
% only points back to the exercises, is left out; the two on checking and counting gates are kept.

\answerchapter{c4:ch}

Varje uttryck nedan är kontrollerat mot sin sanningstabell, rad för rad. Gör samma sak med dina
egna svar: det är den enda kontroll som säger att en gruppering är rätt och inte bara ser rätt ut.

Grindantal räknas med grindar som de ritas i nätet. Där kretsantal efterfrågas räknas med
74-seriens grindar, som har två ingångar (en AND med tre ingångar blir då två AND-grindar i rad).

\answer{c4:ex:1}{Graykod och grannar}
\pt{a}Två rutor som ligger bredvid varandra ska skilja sig i \textbf{exakt en} ingång. Då betyder
två ettor bredvid varandra att utgången inte beror på den ingången, och den kan strykas. I
ordningen \code{00, 01, 10, 11} ändras båda bitarna mellan \code{01} och \code{10}, så de två
mittkolumnerna skulle se ut som grannar utan att vara det.

\pt{b}Ruta 5 är \code{0101}. Byt en bit i taget:

\begin{narrowtable}{@{}ccc@{}}
  \thead{Bit som byts} & \thead{Ny kombination} & \thead{Minterm} \\
  \midrule
  A & 1101 & 13 \\
  B & 0001 & 1 \\
  C & 0111 & 7 \\
  D & 0100 & 4 \\
\end{narrowtable}

\noindent Grannarna är \textbf{1, 4, 7 och 13}. I diagrammet ligger de ovanför, till vänster, till
höger och nedanför ruta 5.

\pt{c}Ruta 0 är \code{0000}, och grannarna är 8 (\code{1000}), 4 (\code{0100}), 2 (\code{0010})
och 1 (\code{0001}). \textbf{Ruta 2} ligger i kolumnen \code{CD = 10} längst till höger och
\textbf{ruta 8} på raden \code{AB = 10} längst ned: båda är grannar över en kant.

\pt{d}En grupp om $2^{k}$ rutor är precis de rutor där $k$ ingångar tar alla sina kombinationer
medan resten är fasta. Tre rutor kan inte vara det: två ingångar ger fyra kombinationer, en ingång
ger två. Tre ettor på rad ringas in som två överlappande par.

\needspace{20\baselineskip}
\answer{c4:ex:2}{Två ingångar}
\pt{a) och b}\leavevmode

\bookfigure{l04_sol_2var}{Övning~\ref{c4:ex:2}: kolumnen \code{B = 1} är en grupp.}%
{c4:fig:sol_2var}

\noindent Båda ettorna ligger i kolumnen \code{B = 1}. \code{A} är 0 i den ena och 1 i den andra,
och stryks:

\begin{console}
X = B
\end{console}

\noindent\code{X} beror \textbf{inte alls} på \code{A}. Kretsen är en ledning från \code{B} till
\code{X}, och \code{A} behöver inte ens kopplas in.

\pt{c}Direkt ur tabellen: raderna \code{01} och \code{11} ger \code{X = A'B + AB}. Bryt ut
\code{B}: \code{X = B(A' + A) = B · 1 = B}. Samma svar. Karnaughdiagrammet gör precis den
förenklingen: \code{A' + A = 1} är det som händer när två grannrutor slås ihop.

\answer{c4:ex:3}{Tre ingångar}
\pt{a}Ettorna ligger i rutorna 0, 1, 4 och 5, alltså på raderna \code{AB = 00} och
\code{AB = 10}. På papperet ligger de överst och nederst, men de \textbf{är} grannar: de skiljer
sig bara i \code{A}, och diagrammets översta och nedersta rad hänger ihop runt kanten
(\secref{c4:a5}).

\needspace{20\baselineskip}
\pt{b}\leavevmode

\bookfigure{l04_sol_3var_b}{Övning~\ref{c4:ex:3}: översta och nedersta raden är en grupp över
kanten.}{c4:fig:sol_3var_b}

\noindent En grupp om fyra, över kanten. \code{A} varierar, \code{C} varierar, och \code{B = 0} i
alla fyra rutorna:

\begin{console}
X = B'
\end{console}

\pt{c}Det minimerade nätet är \textbf{en} inverterare. Den direkta avläsningen,
\code{X = A'B'C' + A'B'C + AB'C' + AB'C}, hade krävt fyra AND-grindar med tre ingångar, en OR-grind
med fyra ingångar och tre inverterare, åtta grindar i allt.

\ptnote{Vanligt fel}Att ringa in raderna \code{00} och \code{10} som två separata par, \code{A'B'}
och \code{AB'}. Svaret \code{A'B' + AB'} är korrekt men inte minimalt, och går att förenkla vidare
till \code{B'}. Glöm inte kanterna.

\needspace{14\baselineskip}
\answer{c4:ex:4}{Från mintermer}
\pt{a}Mintermerna 0, 2, 5 och 7 är \code{000}, \code{010}, \code{101} och \code{111}:

\begin{narrowtable}{@{}cc@{}}
  \thead{ABC} & \thead{X} \\
  \midrule
  000 & 1 \\
  001 & 0 \\
  010 & 1 \\
  011 & 0 \\
  100 & 0 \\
  101 & 1 \\
  110 & 0 \\
  111 & 1 \\
\end{narrowtable}

\needspace{20\baselineskip}
\pt{b}\leavevmode

\bookfigure{l04_sol_3var_xnor}{Övning~\ref{c4:ex:4}: grupperna \code{A'C'} och
\code{AC}.}{c4:fig:sol_3var_xnor}

\noindent Rutorna 0 och 2 ligger på raderna \code{00} och \code{01} i kolumnen \code{C = 0}:
\code{A = 0} och \code{C = 0}, \code{B} varierar, termen blir \code{A'C'}. Rutorna 5 och 7 ligger
på raderna \code{10} och \code{11} i kolumnen \code{C = 1}: \code{A = 1} och \code{C = 1}, termen
blir \code{AC}.

\begin{console}
X = A'C' + AC
\end{console}

\noindent Lägg märke till att raderna \code{11} och \code{10} ligger bredvid varandra i diagrammet,
trots att de kommer i omvänd ordning mot vanlig binär räkning. Det är Graykoden igen.

\pt{c}\code{X} är 1 när \code{A} och \code{C} är lika. Det är \textbf{XNOR}:
\code{X = (A xor C)'}, en enda grind. \code{B} påverkar inte utgången.

\needspace{20\baselineskip}
\answer{c4:ex:5}{Två utgångar}
\pt{a) och b}\leavevmode

\bookfigure{l04_sol_x}{Övning~\ref{c4:ex:5}, \code{X}: raderna \code{AB = 01} och \code{AB = 10}
är två grupper om fyra.}{c4:fig:sol_x}

\noindent\code{X} har ettor på hela raden \code{AB = 01} och hela raden \code{AB = 10}. Två grupper
om fyra, där \code{C} och \code{D} varierar över hela raden:

\begin{console}
X = A'B + AB'
\end{console}

\bookfigure{l04_sol_y}{Övning~\ref{c4:ex:5}, \code{Y}: kolumnerna \code{CD = 00} och
\code{CD = 11} är två grupper om fyra.}{c4:fig:sol_y}

\noindent\code{Y} har ettor i hela kolumnen \code{CD = 00} och hela kolumnen \code{CD = 11}:

\begin{console}
Y = C'D' + CD
\end{console}

\noindent\code{X} påverkas inte av \code{C} och \code{D}: dess grupper är hela \textbf{rader}, där
\code{C} och \code{D} tar alla värden. \code{Y} påverkas inte av \code{A} och \code{B}: dess
grupper är hela \textbf{kolumner}. När en grupp sträcker sig över hela diagrammet i en riktning
betyder det att variablerna i den riktningen saknar betydelse.

\pt{c}\code{X = A'B + AB'} är \textbf{XOR} av \code{A} och \code{B}. \code{Y = C'D' + CD} är
\textbf{XNOR} av \code{C} och \code{D}. Två grindar i stället för fyra AND, två OR och fyra
inverterare.

\needspace{20\baselineskip}
\answer{c4:ex:6}{Hörnen}
\pt{a}\leavevmode

\bookfigure{l04_sol_corners}{Övning~\ref{c4:ex:6}: de fyra hörnen är en grupp, och ruta 15 står
ensam.}{c4:fig:sol_corners}

\noindent Hörnen 0, 2, 8 och 10 är en grupp om fyra, \code{B'D'} (se \secref{c4:a5}). Ruta 15 står
ensam:

\begin{console}
X = B'D' + ABCD
\end{console}

\pt{b}Nej. Ruta 15:s grannar är 7, 11, 13 och 14, och ingen av dem är en etta. En grupp får bara
innehålla ettor, så ruta 15 blir en grupp om en, med alla fyra ingångarna kvar i termen. En ensam
etta är ingen felgruppering, bara en funktion som inte går att förenkla mer just där.

\needspace{20\baselineskip}
\answer{c4:ex:7}{Segment e i en sifferdisplay}
\pt{a) och b}\leavevmode

\bookfigure{l04_sol_segment_e}{Övning~\ref{c4:ex:7}: segment e, med don't care i rutorna 10--15
och två grupper.}{c4:fig:sol_segment_e}

\noindent Ettor i 0, 2, 6 och 8. Utan X-rutorna kan 0, 2 och 8 bara bilda par. Med X i ruta 10 blir
hörnen en grupp om fyra, \code{B'D'}. Ruta 6 går ihop med 2, 14 och 10 i kolumnen \code{CD = 10},
där 14 och 10 är X: \code{CD'}.

\begin{console}
e = B'D' + CD'
\end{console}

\needspace{14\baselineskip}
\pt{c}\leavevmode

\begin{narrowtable}{@{}cccccl@{}}
  \thead{Siffra} & \thead{ABCD} & \thead{\code{B'D'}} & \thead{\code{CD'}} & \thead{e} &
    \thead{Ska vara} \\
  \midrule
  0 & 0000 & 1 & 0 & 1 & tänd \\
  1 & 0001 & 0 & 0 & 0 & släckt \\
  2 & 0010 & 1 & 1 & 1 & tänd \\
  3 & 0011 & 0 & 0 & 0 & släckt \\
  4 & 0100 & 0 & 0 & 0 & släckt \\
  5 & 0101 & 0 & 0 & 0 & släckt \\
  6 & 0110 & 0 & 1 & 1 & tänd \\
  7 & 0111 & 0 & 0 & 0 & släckt \\
  8 & 1000 & 1 & 0 & 1 & tänd \\
  9 & 1001 & 0 & 0 & 0 & släckt \\
\end{narrowtable}

\noindent Alla tio stämmer.

\pt{d}\code{1010} ger \code{B'D' = 1}, så segment e lyser. Displayen visar något som inte är en
siffra. Det är inget problem så länge kretsen framför bara skickar BCD, och det var just det
antagandet som gav oss rätt att använda X-rutorna. Om kombinationen \textbf{kan} komma, är den inte
\enquote{don't care}, och då måste diagrammet ha nollor där.

\answer{c4:ex:8}{Fläkten}
\pt{a) och b}\code{F = 1} när \code{(T = 1 och W = 0) eller C = 1}, och \code{B = 0}. Ettorna blir
mintermerna 2, 6, 8, 10 och 14 (\code{TWCB} = 0010, 0110, 1000, 1010, 1110):

\bookfigure{l04_sol_fan}{Övning~\ref{c4:ex:8}: grupperna \code{CB'} och \code{TW'B'}, den senare
över kanten.}{c4:fig:sol_fan}

\noindent Kolumnen \code{CB = 10} är full: \code{C = 1} och \code{B = 0}, termen \code{CB'}. Raden
\code{TW = 10} har ettor i kolumnerna \code{00} och \code{10}, som är grannar över kanten:
\code{T = 1}, \code{W = 0}, \code{B = 0}, termen \code{TW'B'}.

\begin{console}
F = TW'B' + CB'
\end{console}

\noindent Läs det tillbaka mot specifikationen: \enquote{varmt och fönstret stängt, och inget
brandlarm; eller hög koldioxidhalt, och inget brandlarm}. Det är precis det som stod.

\ptnote{Vanligt fel}Att glömma brandlarmet i en av termerna, till exempel \code{F = TW'B' + C}. Då
går fläkten när koldioxidhalten är hög, även under en brand. En sanningstabell som är härledd ur
kravet, och inte ur en skiss av nätet, fångar det direkt.

\pt{c}Två inverterare (\code{W'} och \code{B'}), en AND-grind med tre ingångar för \code{TW'B'}, en
AND med två för \code{CB'} och en OR med två, fem grindar. Med tvåingångsgrindar blir det sex,
eftersom AND-grinden med tre ingångar blir två (se d). Bryter du ut \code{B'},
\code{F = B'(TW' + C)}, räcker två AND-grindar med två ingångar, en OR och samma två inverterare:
fem grindar, ingen med mer än två ingångar.

\pt{d}Med \code{F = TW'B' + CB'} och tvåingångsgrindar: två inverterare ur \textbf{en 74HC04}, tre
AND-grindar (två för \code{TW'B'}, en för \code{CB'}) ur \textbf{en 74HC08}, och en OR-grind ur
\textbf{en 74HC32}. Tre kretsar, med fyra inverterare, en AND och tre OR över till annat.

\answer{c4:ex:9}{Enbart NAND}
\pt{a}En inverterare för \code{B'}, en AND-grind för \code{AB'} och en OR-grind som lägger till
\code{C}.

\pt{b}Varje AND blir en NAND, OR-grinden blir en NAND, och \code{C}, som går direkt till
OR-grinden, måste inverteras. Inverterare görs av NAND med ingångarna ihopkopplade:

\begin{narrowtable}{@{}ccc@{}}
  \thead{Grind} & \thead{Ingångar} & \thead{Utgång} \\
  \midrule
  N1 & B, B & \code{B'} \\
  N2 & A, N1 & \code{(AB')'} \\
  N3 & C, C & \code{C'} \\
  N4 & N2, N3 & X \\
\end{narrowtable}

\pt{c}\code{X = ((AB')' · C')'}. Med De Morgan, \code{(PQ)' = P' + Q'}:
\code{X = (AB')'' + C'' = AB' + C}. Stämmer.

\pt{d}Fyra NAND-grindar, precis \textbf{en 74HC00}. AND-OR-NOT-versionen hade krävt tre kretsar:
en 74HC04, en 74HC08 och en 74HC32.

\answer{c4:ex:10}{Enbart NOR}
\pt{a}\code{(A + B)C} är en \textbf{produkt av summor} (PS-form): en OR-term \code{(A + B)} och en
term med en enda variabel, \code{C}, ihop med AND.

\pt{b}Varje OR blir en NOR, AND-grinden blir en NOR, och \code{C}, som går direkt till
AND-grinden, inverteras med en NOR med ingångarna ihopkopplade:

\begin{narrowtable}{@{}ccc@{}}
  \thead{Grind} & \thead{Ingångar} & \thead{Utgång} \\
  \midrule
  N1 & A, B & \code{(A + B)'} \\
  N2 & C, C & \code{C'} \\
  N3 & N1, N2 & X \\
\end{narrowtable}

\noindent\code{X = ((A + B)' + C')'}, och med De Morgan, \code{(P + Q)' = P'Q'}:
\code{X = (A + B)'' · C'' = (A + B)C}.

\needspace{12\baselineskip}
\pt{c}\leavevmode

\begin{narrowtable}{@{}ccccccc@{}}
  \thead{A} & \thead{B} & \thead{C} & \thead{N1} & \thead{N2} & \thead{X} &
    \thead{\code{(A + B)C}} \\
  \midrule
  0 & 0 & 0 & 1 & 1 & 0 & 0 \\
  0 & 0 & 1 & 1 & 0 & 0 & 0 \\
  0 & 1 & 0 & 0 & 1 & 0 & 0 \\
  0 & 1 & 1 & 0 & 0 & 1 & 1 \\
  1 & 0 & 0 & 0 & 1 & 0 & 0 \\
  1 & 0 & 1 & 0 & 0 & 1 & 1 \\
  1 & 1 & 0 & 0 & 1 & 0 & 0 \\
  1 & 1 & 1 & 0 & 0 & 1 & 1 \\
\end{narrowtable}

\noindent Tre NOR-grindar, en 74HC02.

\answer{c4:ex:11}{Kretsar och ben}
\pt{a}74HC00, grind 3: ingångar på \textbf{ben 9 och 10}, utgång på \textbf{ben 8}. 74HC02,
grind 3: ingångar på \textbf{ben 8 och 9}, utgång på \textbf{ben 10}. Samma grindnummer, olika ben:
74HC02 har utgången först i varje grupp.

\pt{b}Sex AND-grindar kräver \textbf{två 74HC08} (åtta grindar, två över), och inverteraren
\textbf{en 74HC04}. Tre kretsar.

\pt{c}\leavevmode
\[
  R = \frac{5 - 2,1}{0,005} = 580\ \Omega
\]
Närmaste standardvärde uppåt är \textbf{680~Ω}, som ger \code{I = 2,9 / 680 ≈ 4,3 mA}. Det är gott
nog: en lysdiod lyser tydligt vid några milliampere. 560~Ω, närmast nedåt, ger 5,2~mA och går också
bra.

\pt{d}En CMOS-ingång som inte är ansluten har ingen bestämd nivå (\secref{c4:b4}). Den kan läsas
som 0 eller 1, växla av störningar, eller hamna mitt emellan, där kretsen drar onödigt mycket ström
och kan bli varm. Att det \enquote{fungerar} betyder bara att de oanvända grindarnas utgångar inte
är kopplade till något du tittar på. Men störningen påverkar hela kretsen genom den gemensamma
matningen, och felet kommer och går med handens närhet och luftfuktigheten. Koppla oanvända
ingångar till GND.

\answer{c4:ex:12}{Kontroll: minimerat och ominimerat nät}
\pt{a}\textbf{För hand.} Mintermerna 1, 3, 5, 7, 9 och 11 är \code{0001}, \code{0011},
\code{0101}, \code{0111}, \code{1001} och \code{1011}:

\begin{console}
X = A'B'C'D + A'B'CD + A'BC'D + A'BCD + AB'C'D + AB'CD
\end{console}

\bookfigure{l04_sol_check}{Övning~\ref{c4:ex:12}: grupperna \code{A'D} och \code{B'D}, den senare
över kanten.}{c4:fig:sol_check}

\noindent Alla ettor ligger i kolumnerna \code{CD = 01} och \code{CD = 11}, där \code{D = 1}.
Raderna \code{00} och \code{01} ger en grupp om fyra, \code{A'D}. Raderna \code{00} och \code{10}
ger en grupp om fyra över kanten, \code{B'D}.

\begin{console}
X = A'D + B'D
\end{console}

\pt{b}\textbf{Förutsäg.} Den direkta avläsningen: sex AND med fyra ingångar, en OR med sex, och tre
inverterare (\code{A'}, \code{B'} och \code{C'}; \code{D'} behövs aldrig, eftersom \code{D} står
oinverterad i alla sex termerna). Med tvåingångsgrindar blir varje fyraingångs-AND tre grindar,
alltså 18 AND, och OR-grinden fem, totalt 26 grindar: fem 74HC08, två 74HC32 och en 74HC04,
\textbf{åtta kretsar}.

Det minimerade uttrycket: två inverterare, två AND och en OR, fem grindar i \textbf{tre kretsar}
(74HC04, 74HC08, 74HC32).

\pt{c) och d}De två utgångarna ska vara lika på alla sexton raderna: 1 för raderna 1, 3, 5, 7, 9
och 11, och 0 för övriga. De två uttrycken beskriver samma funktion, så varje avvikelse är ett fel i
ett av näten, oftast en felkopplad ledning i CircuitVerse, inte i minimeringen. Om de skiljer sig
på en rad: kontrollera den raden i båda uttrycken för hand. Det uttryck som ger fel värde där är det
felaktiga, och den AND-term eller grupp som borde ha täckt raden är den som ska granskas.

\pt{e}\code{X = D(A' + B') = D · (AB)'}: en NAND-grind och en AND-grind, i en 74HC00 och en
74HC08, \textbf{två kretsar}. Eller med enbart NAND: \code{N1 = (AB)'}, \code{N2 = (N1 · D)'},
\code{X = (N2 · N2)' = N1 · D}, tre NAND-grindar i \textbf{en} 74HC00. Från åtta kretsar till en,
för samma sanningstabell. Det är hela poängen med att minimera, och med att ibland tänka ett steg
längre än diagrammet.

\needspace{14\baselineskip}
\answer{c4:ex:13}{Multiplexern}
\pt{a}\leavevmode

\begin{narrowtable}{@{}cccc@{}}
  \thead{S} & \thead{A} & \thead{B} & \thead{X} \\
  \midrule
  0 & 0 & 0 & 0 \\
  0 & 0 & 1 & 1 \\
  0 & 1 & 0 & 0 \\
  0 & 1 & 1 & 1 \\
  1 & 0 & 0 & 0 \\
  1 & 0 & 1 & 0 \\
  1 & 1 & 0 & 1 \\
  1 & 1 & 1 & 1 \\
\end{narrowtable}

\noindent När \code{S = 0} är \code{X} en kopia av \code{B}, när \code{S = 1} en kopia av \code{A}.

\needspace{20\baselineskip}
\pt{b}\leavevmode

\bookfigure{l04_sol_mux}{Övning~\ref{c4:ex:13}: grupperna \code{S'B} och
\code{SA}.}{c4:fig:sol_mux}

\begin{console}
X = S'B + SA
\end{console}

\pt{c}En inverterare för \code{S'}, två AND-grindar och en OR-grind. Den övre AND-grinden får
\code{S'} och \code{B}, den nedre \code{S} och \code{A}. Eftersom \code{S} och \code{S'} aldrig är
1 samtidigt kan högst en av AND-grindarna släppa igenom något: \code{S} fungerar som en
strömbrytare som öppnar den ena vägen och stänger den andra. Samma idé, med fler väljarbitar, sitter
inuti mikrokontrollern och väljer vilken av dess analoga ingångar som ska mätas.
